\begin{textsample}{NEG Dim 3 – 2021\_c5 – Score 2.00 – 2021\_c5/enade\_2021\_pedag\_7adb919c15.txt}  \label{ex:f3_neg_196}
CAPÍTULO 5 DISTRIBUIÇÃO DOS CONCEITOS 5.1 PANORAMA NACIONAL DA DISTRIBUIÇÃO DOS CONCEITOS O Conceito Enade é um indicador de qualidade que avalia os cursos por intermédio dos desempenhos dos estudantes a partir dos resultados do Enade.
Os valores possíveis vão de 1 ( pior situação ) a 5 ( melhor situação ).
Os cursos com apenas um ou sem nenhum concluinte participante não obtêm o Conceito Enade, ficando Sem Conceito ( SC ).
Na Tabela 5.1, são apresentadas a quantidade e a distribuição de cursos de Pedagogia participantes do Enade/2021, por faixa de conceito e Grande Região.
Enfatiza -se, mais uma vez, que a diferença entre os cursos tabulados neste capítulo e no capítulo 2 corresponde aos cursos que não foram avaliados, ou seja, àqueles, em princípio, sem alunos concluintes que tivessem sido inscritos para a prova.
Observando -se os dados da Tabela 5.1, nota -se que, dos 1.249 cursos participantes, 551 ( 44,1 \% ) classificaram -se com conceito 3, o valor modal.
Esse foi também o conceito modal na maioria das regiões : Nordeste ( 44,4 \% ), Sudeste ( 44,5 \% ), Sul ( 44,2 \% ) e Centro-Oeste ( 45,3 \% ).
O conceito 2 foi o segundo mais frequente em nível nacional ( 28,7 \%, correspondendo a 358 cursos ), tendo sido o conceito modal na região Norte ( 42,0 \% ).
O conceito 4 foi o terceiro ( 18,6 \%, correspondendo a 232 cursos ).
Houve, ainda, 50 cursos ( 4,0 \% ) que receberam o conceito 5, e 44 cursos ( 3,5 \% ) que receberam o conceito 1.
Dos 1.249 cursos de Pedagogia, 14 ( 1,1 \% ) ficaram Sem Conceito ( SC ).
A região Norte participou com 69 cursos, o correspondente a 5,5 \% do total nacional.
Desses, 29 cursos receberam o conceito 2, o que equivale a 42,0 \% do total regional, o conceito modal na região, como já comentado.
A cinco cursos ( 7,2 \% ) atribuiu -se o conceito 1, e a 26 cursos ( 37,7 \% ), o conceito 3.
Nenhum curso ficou Sem Conceito ( SC ), um curso ( 1,4 \% ) recebeu o conceito 5, e oito cursos ( 11,6 \% ), o conceito 4, nessa região.
A região Nordeste participou com 239 cursos, o correspondente a 19,1 \% do total nacional.
Nessa região, 106 cursos ( 44,4 \% em termos regionais ) obtiveram conceito 3, o conceito modal para a região, como já comentado.
Os conceitos 1 e 2 foram atribuídos, respectivamente, a 12 e a 69 cursos ( 5,0 \% e 28,9 \%, respectivamente ).
Ainda nessa região, 44 cursos ( 18,4 \% ) foram avaliados com o conceito 4, e sete cursos ( 2,9 \% ) foram avaliados com o conceito 5.
Um curso ficou Sem Conceito ( SC ).
Para os 564 cursos participantes da região Sudeste ( 45,2 \% do total nacional ), o conceito 3 foi o conceito modal, como já anteriormente assinalado, tendo sido obtido por 251 cursos ( 44,5 \% ).
O conceito 1 foi atribuído a dez cursos ( 1,8 \% ), e o conceito 2, a 154 cursos ( 27,3 \% ).
Receberam os conceitos 4 e 5, respectivamente, 110 cursos ( 19,5 \% ) e 31 cursos ( 5,5 \% ).
Nessa região, oito cursos ficaram Sem Conceito ( SC ).
Os 240 cursos da região Sul corresponderam a 19,2 \% do total nacional.
Como já indicado, houve predominância do conceito 3 ( conceito modal ), atribuído a 44,2 \% dos cursos da região, ou seja, atribuído a 106 dos 240 cursos participantes na região Sul.
O conceito 1 foi atribuído a quatros cursos ( 1,7 \% ) ; o conceito 2, a 59 cursos ( 24,6 \% ) ; e os conceitos 4 e 5, a 58 ( 24,2 \% ) e a oito cursos ( 3,3 \% ), respectivamente.
Na região Sul, cinco cursos ficaram Sem Conceito ( SC ).
Já dos 137 cursos participantes na região Centro-Oeste ( 11,0 \% do total nacional ), 62 ( 45,3 \% em termos regionais ) receberam conceito 3, o conceito modal, como já destacado.
Treze cursos ( 9,5 \% ) obtiveram o conceito 1, 47 cursos ( 34,3 \% ) obtiveram o conceito 2, 12 ( 8,8 \% ), o conceito 4, e outros três ( 2,2 \% ), o conceito 5.
Nessa região, nenhum curso ficou Sem Conceito ( SC ).
No Gráfico 5.1, é apresentada a distribuição cumulativa do Conceito Enade, segundo a Grande Região.
Em um gráfico de distribuição cumulativa, em cada valor de abscissa ( eixo x ), o valor apresentado no eixo das ordenadas ( eixo y ) é o valor acumulado da variável ( Conceito Enade ).
Observando -se a linha vertical pontilhada, que ajuda a visualizar o percentual de conceitos acumulados até o valor 3, ou seja, os conceitos 1, 2 e 3, pode -se notar que, por exemplo, a região Centro-Oeste ( linha laranja ) apresentou quase 90 \% dos cursos com o conceito 3 ou abaixo desse valor.
A região Sul ( linha rosa ), por outro lado, apresentou menos de 72 \% dos cursos na mesma situação : com o conceito 3 ou abaixo.
Em linhas gerais, em um gráfico de distribuição cumulativa, poligonais mais à esquerda representam as Grandes Regiões com a pior distribuição, e poligonais mais à direita, as Grandes Regiões com os melhores conceitos.
A região Norte ( linha azul ) e a Centro-Oeste ( linha laranja ) ( linha laranja ) apresentaram os cursos com os conceitos mais baixos, ao passo que a região Sul ( linha rosa ) e a Sudeste ( linha vermelha ) apresentaram a distribuição com os valores maiores.
Cursos em IES situadas nas regiões Nordeste ( linha verde ) apresentaram uma situação intermediária. 5.2 CONCEITOS POR CATEGORIA ADMINISTRATIVA, POR MODALIDADE DE ENSINO E GRANDE REGIÃO Na Tabela 5.2, é apresentada a distribuição dos cursos de Pedagogia participantes do Enade/2021, por Categoria Administrativa e Modalidade de Ensino, de acordo com os conceitos por eles alcançados, segundo as Grandes Regiões brasileiras.
Dos 1.249 cursos participantes, 334 ( 26,7 \% ) eram ministrados em Instituições Públicas, e 915 ( 73,3 \% ), em Instituições Privadas.
Quanto à Modalidade de Ensino, 190 cursos eram a Distância, e 1.059, Presenciais.
De acordo com as informações da Tabela 5.2, em termos nacionais, entre Instituições Públicas e Privadas, observa -se que, dos 50 cursos avaliados com conceito 5, 29 foram oferecidos em IES Públicas, e 21, em IES Privadas.
Dos 334 cursos participantes de IES Públicas, o conceito 3 foi o valor modal, atribuído a 140 cursos ( 41,9 \% ).
Dos demais cursos de Instituições Públicas participantes, cinco cursos ( 1,5 \% ) foram avaliados com o conceito 1, 56 cursos ( 16,8 \% ) foram avaliados com o conceito 2, 102 cursos ( 30,5 \% ) foram avaliados com o conceito 4 e, 29 cursos ( 8,7 \% ), com o conceito 5.
Nessa categoria, dois cursos ( 0,6 \% ) ficaram Sem Conceito ( SC ).
Na rede privada, o conceito modal foi o 3, atribuído a 411 cursos, o correspondente a 44,9 \% dos 915 cursos da categoria.
Dos demais cursos participantes, 39 ( 4,3 \% ) receberam o conceito 1, e 302 ( 33,0 \% ), o conceito 2.
O conceito 4 foi atribuído a 130 cursos ( 14,2 \% ), e o conceito 5, a 21 cursos ( 2,3 \% ).
Nessa Categoria Administrativa, 12 cursos ( 1,3 \% ) ficaram Sem Conceito ( SC ).
Dos 1.059 cursos participantes da Modalidade Presencial, o conceito 3 foi o valor modal, atribuído a 478 cursos ( 45,1 \% ).
Dos demais cursos nessa Modalidade de Ensino, 31 cursos ( 2,9 \% ) receberam o conceito 1 ; 288 cursos ( 27,2 \% ) receberam o conceito 2 ; 209 cursos ( 19,7 \% ), o conceito 4 ; e 44 cursos ( 4,2 \% ), o conceito 5.
Nessa Modalidade de Ensino, nove cursos ( 0,8 \% ) ficaram Sem Conceito ( SC ).
Dos 190 cursos participantes da Modalidade a Distância, o conceito 3 foi o valor modal, atribuído a 73 cursos ( 38,4 \% ).
Dos demais cursos nessa Modalidade de Ensino, 13 cursos ( 6,8 \% ) receberam o conceito 1 ; 70 cursos ( 36,8 \% ) receberam o conceito 2 ; 23 cursos, ( 12,1 \% ), o conceito 4 ; e 6 cursos ( 3,2 \% ), o conceito 5.
Nessa Modalidade de Ensino, cinco cursos ( 2,6 \% ) ficaram Sem Conceito ( SC ).
Na análise por região, observa -se que, na região Norte, as Instituições Privadas participaram com 35 cursos ( 50,7 \% do total regional ), e nenhum curso ficou Sem Conceito ( SC ).
O conceito modal para as Instituições Privadas na região foi o 2, atribuído a 18 cursos, correspondendo a 51,4 \% dessa categoria na região.
Quatro cursos foram avaliados com o conceito 1 ; 11, com o conceito 3 ; dois, com conceito 4.
A nenhum curso foi atribuído o conceito 5.
As Instituições Públicas participaram com 34 cursos na região Norte ( 49,3 \% do total regional ), e nenhum curso ficou Sem Conceito ( SC ).
O conceito modal nas Instituições Públicas na região foi o 3, atribuído a 15 cursos, o correspondente a 44,1 \% dessa categoria na região.
Um curso foi avaliado com conceito 1 ; 11 receberam o conceito 2 ; seis receberam o conceito 4 ; e apenas um curso obteve o conceito 5.
Dos cursos oferecidos na região Norte, 63 eram presenciais, e nenhum ficou Sem conceito ( SC ).
Dos 63 cursos oferecidos, cinco obtiveram o conceito 1 ; 26, o conceito 2 ( modal ) ; 24, o conceito 3 ; sete, o conceito 4 ; e um, o conceito 5.
Seis eram na Modalidade a Distância, tendo três recebido o conceito modal 2 ; dois, o conceito 3 ; e um, o conceito 4.
Nenhum curso ficou Sem Conceito ( SC ) ou recebeu os demais conceitos.
Na região Nordeste, a rede privada concentrou 136 dos 239 cursos participantes, o equivalente a 56,9 \% do total da região.
Entre os cursos oferecidos por IES Privadas no Nordeste, 57 cursos foram avaliados com conceito 2, o conceito modal.
Doze cursos obtiveram o conceito 1 ; 54, o conceito 3 ; 11, o conceito 4 ; e apenas um curso obteve o conceito 5.
Nessa Categoria Administrativa e Grande Região, um curso ficou Sem Conceito ( SC ).
As Instituições Públicas da região Nordeste participaram com 103 cursos ( 43,1 \% do total da região ).
Desses, 52 ( 50,5 \% ) obtiveram conceito 3, o valor modal, seguidos por 33 cursos, que receberam o conceito 4 ; 12, o conceito 2 ; e seis, o conceito 5.
Nenhum curso ficou Sem Conceito ( SC ) ou com conceito 1 nessa região.
Dos cursos oferecidos na região Nordeste, 217 eram Presenciais : um ficou Sem conceito ( SC ), oito ficaram com o conceito 1 ; 62, com o conceito 2 ; 96, com o conceito 3, o valor modal ; 43, com o conceito 4 ; e sete, com o conceito 5.
Na Modalidade a Distância, eram 22 cursos, tendo quatro obtido o conceito 1 ; sete, o conceito 2 ; dez, o conceito 3, o valor modal ; um, o conceito 4.
Nessa região, nenhum curso ficou Sem Conceito ( SC ) ou recebeu o conceito 5.
Na região Sudeste, a proporção de cursos da rede privada, 83,9 \%, correspondeu a 473 dos 564 cursos participantes.
Nessa combinação de Categoria Administrativa e Grande Região, o conceito modal foi o 3, atribuído a 223 cursos.
Seis cursos ficaram Sem Conceito ( SC ).
Os demais foram avaliados com os seguintes conceitos : 1 ( dez cursos ), 2 ( 143 cursos ), 4 ( 76 cursos ) e 5 ( 15 cursos ).
Dos 91 cursos oferecidos por Instituições Públicas na região Sudeste, o conceito modal foi o 4, atribuído a 34 cursos.
Dos demais cursos, 11 receberam o conceito 2 ; 28 receberam o conceito 3 ; e 16, o conceito 5.
Nessa combinação de Categoria Administrativa e Grande Região, dois cursos ficaram Sem Conceito ( SC ), e nenhum recebeu o conceito 1.
Também na região Sudeste, grande parcela dos cursos foi oferecida na Modalidade Presencial : 471 de 564 cursos.
Para esses, a moda foi o conceito 3, com 215 cursos.
Os demais cursos Presenciais da região foram alocados nos conceitos 1 ( nove cursos ), 2 ( 122 cursos ), 4 ( 97 cursos ) e 5 ( 25 cursos ).
Ficaram Sem Conceito ( SC ) três cursos Presenciais na região Sudeste.
Na região Sudeste, 93 cursos foram oferecidos na Modalidade Educação a Distância, tendo sido um curso com conceito 1 ; 32 com conceito 2 ; 36 com conceito 3, o valor modal ; 13 com o conceito 4 ; e seis com o conceito 5.
Na região Sudeste, cinco cursos a Distância ficaram Sem Conceito ( SC ).
As Instituições Privadas concentraram 188 dos 240 cursos participantes da região Sul, 78,3 \% do total regional.
Desses, 85 cursos obtiveram conceito 3, o conceito modal.
Nessa combinação de Categoria Administrativa e Grande Região, quatro cursos receberam o conceito 1 ; 53 receberam o conceito 2 ; outros 37, o conceito 4 ; e ainda quatro, o conceito 5.
Cinco cursos ficaram Sem Conceito ( SC ).
As Instituições Públicas na região Sul participaram com 52 cursos ( 21,7 \% dos cursos da região ), aos quais foram atribuídos os conceitos 2 ( 6 cursos ), 3 e 4 ( 21 cursos cada, modais ) e 5 ( quatro cursos ).
Nenhum curso da região Sul oferecido por IES Públicas ficou Sem Conceito ( SC ) ou recebeu o conceito 1.
Na região Sul, 186 dos cursos foram oferecidos na Modalidade Presencial.
Para esses, a moda foi o conceito 3, com 85 cursos.
Os demais cursos Presenciais da região foram alocados nos conceitos 1 ( um curso ), 2 ( 36 cursos ), 4 ( 51 cursos ) e 5 ( oito cursos ).
Ficaram Sem Conceito ( SC ) cinco cursos.
Na Modalidade Educação a Distância, foram oferecidos 54 cursos, sendo três cursos com conceito 1 ; 23 com conceito 2, o valor modal ; 21 com conceito 3 ; e sete com o conceito 4.
Nenhum curso ficou Sem Conceito ( SC ) ou com o conceito 5 nessa região.
Na região Centro Oeste, 83 dos 137 cursos participantes eram de Instituições Privadas ( 60,6 \% em termos regionais ).
Desses, 38 cursos receberam conceito 3, o valor modal.
Dos demais cursos, nove receberam o conceito 1 ; 31 receberam o conceito 2 ; quatro, o conceito 4 ; e um curso recebeu o conceito 5.
Nenhum curso ficou Sem Conceito ( SC ).
Quanto aos 54 cursos oferecidos por Instituições Públicas na região Centro Oeste ( 39,4 \% do total regional ), 24 cursos apresentaram o conceito 3, conceito modal ; quatro cursos obtiveram o conceito 1 ; 16, o conceito 2 ; oito, o conceito 4 ; e dois cursos, o conceito 5.
Nenhum curso ficou Sem Conceito ( SC ) nessa região.
Também na região Centro Oeste, a maioria dos cursos foi Presencial : 122 ( 89,1 \% ).
Para esses, a moda foi o conceito 3, com 58 cursos ( 47,5 \% ).
Os demais cursos Presenciais da região foram alocados nos conceitos 1 ( oito cursos, 6,6 \% ), 2 ( 42 cursos, 34,4 \% ), 4 ( 11 cursos, 9,0 \% ) e 5 ( três cursos, 2,5 \% ).
Nenhum curso Presencial da região ficou Sem Conceito ( SC ).
Para os 15 cursos a Distância ( 10,9 \% ), a moda bimodal foram os conceitos 1 e 2 ( 5 cursos cada ).
Os demais cursos a Distância da região receberam os conceitos 3 ( quatro cursos ) e 4 ( um curso ).
Nenhum curso a Distância da região ficou Sem Conceito ( SC ) ou com o conceito 5.
No Gráfico 5.2, apresenta -se a distribuição do Conceito Enade, segundo a Categoria Administrativa da IES.
Os cursos em IES Públicas ( linha azul ) apresentam uma poligonal mais à direita do que os em IES Privadas ( linha verde ) e, consequentemente, uma distribuição de Conceito Enade com valores maiores. 5.3 CONCEITOS POR ORGANIZAÇÃO ACADÊMICA E GRANDE REGIÃO Na Tabela 5.3, encontra -se a distribuição dos conceitos atribuídos aos cursos participantes do Enade/2021, na Área de Pedagogia, por Organização Acadêmica, segundo as Grandes Regiões brasileiras.
Dos 1.249 cursos de Pedagogia participantes, 488 eram oferecidos em Universidades, 277 em Centros Universitários, 456 em Faculdades e 19 em CEFET/IF.
Esta distribuição corresponde a, respectivamente, 39,1 \%, 22,2 \%, 37,2 \% e 1,5 \% dos cursos.
De acordo com os dados apresentados, de todos os 50 cursos avaliados com o conceito 5, 40 eram vinculados a Universidades.
Os cursos vinculados a esse tipo de Organização Acadêmica tiveram o conceito 3 como conceito modal, atribuído a 223 cursos ( 45,7 \% ).
Os demais cursos vinculados a Universidades receberam os conceitos 1 ( quatro cursos ), 2 ( 89 cursos ), 4 ( 127 cursos ) e conceito 5 ( 40 cursos, como já mencionado ).
Cinco cursos ficaram Sem Conceito ( SC ).
Dos cursos em Centros Universitários, o conceito modal também foi o 3, atribuído a 131 cursos ( 47,3 \% ).
Cinco cursos vinculados a esse tipo de Organização Acadêmica ficaram Sem Conceito ( SC ), e os outros cursos receberam os conceitos 1 ( 12 cursos ), 2 ( 76 cursos ), 4 ( 49 cursos ) e 5 ( quatro cursos ).
Dos 465 cursos mantidos por Faculdades, 192 ( 41,3 \% ) receberam conceito 3, o conceito modal.
O segundo conceito mais frequentemente atribuído aos cursos oferecidos por essa categoria de IES foi o conceito 2, recebido por 189 cursos.
Na sequência, 46 cursos obtiveram conceito 4 ; 28 cursos, o conceito 1 ; e seis cursos, o conceito 5.
Dos cursos oferecidos em Faculdades, quatro ficaram Sem Conceito ( SC ).
Dos 19 cursos mantidos por CEFET/IF, dez ( 52,6 \% ) receberam conceito 4, o conceito modal.
O segundo conceito mais frequentemente atribuído aos cursos oferecidos por essa categoria de IES foi o conceito 3, recebido por cinco cursos.
Na sequência, quatro cursos obtiveram o conceito 2, e nenhum curso ficou Sem Conceito ( SC ) ou com os conceitos 1 e 5.
Considerando -se, separadamente, as regiões brasileiras, verifica -se que, na região Norte, as Universidades participaram com 36 cursos avaliados, aos quais foram atribuídos os seguintes conceitos : 2, a 12 cursos ; 3, o conceito modal, a 15 cursos ; 4 a sete cursos ; 1, a um curso ; e 5, também a um curso.
Nenhum curso ficou Sem Conceito ( SC ) nessa combinação de categoria e região.
Os Centros Universitários da região Norte foram representados por dez cursos, os quais sete obtiveram conceito 2, o valor modal, e três cursos, o conceito 3.
Os conceitos 1, 4 e 5 não foram atribuídos a nenhum curso.
Do mesmo modo, nenhum curso ficou Sem Conceito ( SC ).
As Faculdades participaram com 23 cursos na região Norte, os quais receberam os conceitos 1 ( quatro cursos ), 2 ( dez cursos, conceito modal ), 3 ( oito cursos ) e 4 ( um curso ).
Nenhum curso ficou Sem Conceito ( SC ), e nenhum curso oferecido em Faculdades da região obteve o conceito 5.
Nessa região e Organização Acadêmica, nenhum curso foi oferecido em CEFET/IF.
Na região Nordeste, as Universidades participaram com 104 dos 239 cursos da Área de Pedagogia oferecidos.
O conceito modal foi 3, atribuído a 57 cursos.
Os 33 cursos obtiveram o conceito 4 ; seis cursos, o conceito 5 ; e outros oito, o conceito 2.
Nenhum curso oferecido por Universidades no Nordeste recebeu o conceito 1 ou ficou Sem Conceito ( SC ).
Os Centros Universitários contaram com 38 cursos participantes na região Nordeste, 19 dos quais receberam o conceito modal 3.
Para os demais cursos, foram atribuídos os conceitos 1 ( quatro cursos ), 2 ( nove cursos ) e 4 ( cinco cursos ).
Um curso oferecido por Centros Universitários no Nordeste ficou Sem Conceito ( SC ) ; nenhum recebeu o conceito 5.
As Faculdades foram representadas por 97 cursos na região Nordeste, 52 dos quais receberam conceito 2, o conceito modal.
Aos demais cursos foram atribuídos os conceitos 1 ( oito cursos ), 3 ( 30 cursos ), 4 ( seis cursos ) e 5 ( um curso ).
Nenhum curso ficou Sem Conceito ( SC ).
Nessa combinação de região e Organização Acadêmica, nenhum curso foi oferecido em CEFET/IF.
Na região Sudeste, as Universidades concentraram 188 dos 564 cursos de Pedagogia da região.
Dos cursos oferecidos em Universidades, nessa região, o conceito modal foi 3, atribuído a 79 cursos.
Os demais cursos receberam os conceitos 2 ( 39 cursos ), 4 ( 43 cursos ) e 5 ( 25 cursos ).
Dois cursos ficaram Sem Conceito ( SC ), e nenhum recebeu o conceito 1.
Os Centros Universitários participaram com 147 cursos na região Sudeste, dos quais 69 obtiveram o conceito modal 3.
O conceito 1 foi atribuído a dois cursos ; o conceito 2, a 37 cursos ; o conceito 4, a 34 cursos ; e o conceito 5, a dois cursos.
Três cursos ficaram Sem Conceito ( SC ) nessa combinação de Organização Acadêmica e região.
Já as Faculdades foram representadas por 221 cursos na região Sudeste, dos quais três ficaram Sem Conceito ( SC ).
O conceito modal 3 foi atribuído a 101 cursos.
Aos demais cursos foram atribuídos os conceitos 1 ( oito cursos ), 2 ( 77 cursos ), 4 ( 28 cursos ) e 5 ( quatro cursos ).
Para os oito cursos participantes de CEFET/IF, na região Sudeste, os conceitos foram : 2 ( um curso ), 3 ( dois cursos ) e 4 ( cinco cursos ), o valor modal.
Nenhum curso ficou Sem Conceito ( SC ) ou recebeu os conceitos 1 e 5 nesse tipo de organização.
Dos 240 cursos da região Sul, 103 eram mantidos por Universidades, para os quais o conceito modal foi 3, atribuído a 44 cursos.
Os demais cursos receberam os conceitos 2 ( 14 cursos ), 4 ( 37 cursos ), 5 ( cinco cursos ) ; nenhum curso recebeu o conceito 1.
Nesse tipo de organização, três cursos da região Sul ficaram Sem Conceito ( SC ).
Dos 55 cursos participantes de Centros Universitários da região Sul, um ficou Sem Conceito ( SC ), e 23 cursos receberam conceito 3, o valor modal.
Os conceitos restantes ficaram assim distribuídos : dois cursos ( conceito 1 ) ; 18 cursos ( conceito 2 ) ; nove cursos ( conceito 4 ) ; e dois cursos ( conceito 5 ).
Quanto aos 74 cursos vinculados a Faculdades na região Sul, 37 cursos receberam o conceito 3, valor modal, e somente um curso ficou Sem Conceito ( SC ).
Os demais cursos receberam os conceitos 1 ( dois cursos ), 2 ( 24 cursos ), 4 ( nove cursos ) e 5 ( um curso ).
Dos oito cursos oferecidos em CEFET/IF, dois cursos foram avaliados com o conceito 3.
Foi registrada uma distribuição bimodal, já que três cursos obtiveram o conceito 2 e outros três cursos obtiveram o conceito 4.
Nesse tipo de organização, nenhum curso da região Sul ficou Sem Conceito ( SC ) ou com os conceitos 1 e 5.
Na região Centro-Oeste, 57 dos 137 cursos oferecidos eram mantidos por Universidades.
Para esse tipo de Organização Acadêmica, o conceito modal foi 3, atribuído a 28 cursos ; e, nenhum curso da região ficou Sem Conceito ( SC ).
Os demais cursos receberam os conceitos 1 ( três cursos ), 2 ( 16 cursos ), 4 ( 7 cursos ) e 5 ( três cursos ).
Os Centros Universitários da região Centro-Oeste contaram com 27 cursos : quatro cursos foram avaliados com o conceito 1 ; cinco cursos, com o conceito 2 ; 17 cursos, com o conceito 3 ( modal ) ; um curso recebeu o conceito 4.
Não houve cursos avaliados com o conceito 5 nem cursos Sem Conceito ( SC ).
Dos 50 cursos oferecidos em Faculdades na região Centro-Oeste, seis cursos receberam o conceito 1 ; 26 cursos, o conceito 2 ( modal ) ; 16 cursos, o conceito 3 ; e dois cursos, o conceito 4.
Nessa região, nenhum curso oferecido em Faculdades ficou Sem Conceito ( SC ) ou obteve o conceito 5.
Dos três cursos oferecidos em CEFET/IF, dois cursos foram avaliados com o conceito 4 ( modal ), e um curso obteve o conceito 3.
Nesse tipo de organização, nenhum curso da região Centro-Oeste ficou Sem Conceito ( SC ) ou recebeu os conceitos 1, 2 e 5.
No Gráfico 5.3, apresenta -se a distribuição do Conceito Enade dos cursos de Pedagogia, segundo a Organização Acadêmica da IES.
Nos CEFET/IF ( linha laranja ), todos os cursos têm conceito 4 ou abaixo disso, já que a poligonal cumulativa alcança 100 \% para o conceito 4.
Além disso, os cursos de CEFET/IF apresentam uma poligonal mais à direita do do que as demais, denotando uma melhor distribuição de Conceitos Enade.
No outro extremo, encontram -se os cursos mantidos por Faculdades ( linha vermelha ), apresentando as respectivas poligonais mais à esquerda e, consequentemente, uma distribuição de Conceitos Enade com os menores valores.
Os cursos mantidos por Centros Universitários ( linha verde ) e Universidades ( linha azul ) apresentam poligonais mais centrais e, consequentemente, uma distribuição de Conceitos Enade com os valores intermediários.

% matched lemmas: 
\end{textsample}
